% =====================================================================
% Sensitivity analysis + unified ranking -- methodology paragraph (English).
% Drop-in subsection for the APEN paper. Assumes: graphicx, amsmath.
% (Currency is written "EUR"; replace by \euro / \texteuro if you load eurosym
%  or textcomp.) Figure paths point to results_meso/. Adjust \label / \ref to
%  your numbering. EoL thresholds is used as the illustrative Pareto example;
%  swap sens_eol_pareto.pdf for sens_hthresholds_pareto.pdf if preferred.
% =====================================================================

\subsection{Robustness of the EMS ranking}
\label{sec:sensitivity}

We assess the robustness of the energy-management strategies (EMS) to the
uncertain parameters of the ageing and cost models through Monte-Carlo
sensitivity analyses. For each analysis, the parameters are drawn uniformly
within the intervals specified below and \emph{the same draws are reused for
every EMS}, so that the comparison between strategies is not polluted by sampling
noise. Each draw is propagated through a full 25-year simulation, yielding for
every strategy a two-dimensional outcome: the loss-of-power-supply probability
(LPSP) and the degradation cost.

\paragraph{Sampled parameters.}
Three families of parameters are varied:
\begin{itemize}
  \item \textbf{End-of-life (EoL) thresholds} of the three components: the
        battery EoL is drawn in $[0.6,\,0.8]$ and the PEMFC and PEMWE EoL in
        $[0.85,\,0.95]$ ($N=200$ paired draws per strategy).
  \item \textbf{Degradation-function thresholds of the hydrogen components}
        (PEMFC/PEMWE), i.e.\ the regime thresholds of the piecewise degradation
        model: $f_\mathrm{high}\in[0.70,\,0.90]$, $f_\mathrm{low}\in[0.005,\,0.020]$,
        $f_{30}\in[0.25,\,0.35]$ and $f_{60}\in[0.50,\,0.70]$ ($N=200$ paired
        draws per strategy).
  \item \textbf{Component replacement-cost weights} (the coefficients that weight
        each component's degradation in the cost axis): each weight is varied by
        $\pm30\%$ around its nominal value ($N=5000$, $95\%$ confidence interval).
        Because the LPSP and the replacement schedule are invariant to these
        weights, only the degradation-cost axis dilates, which is handled
        analytically.
\end{itemize}

\paragraph{Robustness to SoH-estimation error.}
The RB2(SoH) strategy acts on the \emph{estimated} state of health (SoH) rather
than on its true value. To check that this dependence is not a liability, we add
a $2\%$ Gaussian noise to the SoH estimate ($N=200$). The resulting performance
remains very close to the exact-information case: the LPSP stays at
$2.66\%$ (range $2.64$--$2.68\%$) against $2.65\%$ with the true SoH, and the
degradation cost at $58.7$~kEUR (range $58.5$--$58.8$~kEUR) against $58.5$~kEUR.
The strategy is therefore essentially insensitive to realistic SoH-estimation
noise.

\paragraph{Unified ranking through a monetary value of the LPSP.}
Each analysis produces a cloud of points in the (LPSP, degradation cost) plane,
which is convenient for inspecting one analysis at a time but does not, by
itself, rank the strategies across all cases. To obtain a single ordering, we
assign a monetary value to the unserved load through a tiered
\emph{value of lost load} (VOLL) and add it to the degradation cost:
\begin{equation}
  C_\mathrm{tot} \;=\; C_\mathrm{deg} \;+\;
  \mathrm{VOLL}(\mathrm{LPSP}) \times E_\mathrm{unserved},
  \qquad
  E_\mathrm{unserved} \;=\; \frac{\mathrm{LPSP}}{100}\,E_\mathrm{ref},
  \label{eq:total_cost}
\end{equation}
where $E_\mathrm{ref}=273.4$~MWh is the net planned load energy over the 25-year
horizon, and the VOLL increases by tiers with the severity of the shortfall:
\begin{equation}
  \mathrm{VOLL}(\mathrm{LPSP}) =
  \begin{cases}
    2~\mathrm{EUR/kWh}  & \text{if } \mathrm{LPSP} < 5\%, \\
    5~\mathrm{EUR/kWh}  & \text{if } 5\% \le \mathrm{LPSP} < 10\%, \\
    10~\mathrm{EUR/kWh} & \text{if } \mathrm{LPSP} \ge 10\%.
  \end{cases}
  \label{eq:voll}
\end{equation}
The single scalar $C_\mathrm{tot}$ in \eqref{eq:total_cost} makes the strategies
directly comparable within and across all sensitivity cases.

\paragraph{Presentation.}
We report the results with two complementary figures. First, as an illustrative
example, a Pareto plane (degradation cost versus LPSP) on which each strategy
carries a $1\sigma$/$2\sigma$ confidence ellipse summarising its Monte-Carlo
dispersion (Fig.~\ref{fig:pareto_example}, here for the EoL-threshold analysis).
Second, a ranking heatmap of $C_\mathrm{tot}$ (strategy $\times$ sensitivity
case, rank~1 = lowest cost), which condenses all analyses into a quick visual
summary (Fig.~\ref{fig:ranking}). Across every case, RB2(SoH) is consistently the
lowest-cost strategy, followed by RB2 and the battery-dominated 100--0 and
75--25 strategies, whereas the strategies that incur a high LPSP (SoC06, 25--75,
0--100) are strongly penalised once the unserved load is priced.

% ---------------------------------------------------------------------
\begin{figure}[t]
  \centering
  \includegraphics[width=0.8\linewidth]{results_meso/sens_eol_pareto.pdf}
  \caption{Illustrative sensitivity analysis (EoL thresholds): Pareto plane of
  the ten EMS, each point at its nominal thresholds with a $1\sigma$/$2\sigma$
  confidence ellipse summarising the Monte-Carlo dispersion over the sampled
  EoL thresholds.}
  \label{fig:pareto_example}
\end{figure}

\begin{figure}[t]
  \centering
  \includegraphics[width=\linewidth]{results_meso/voll_ranking.pdf}
  \caption{Unified ranking of the ten EMS by total cost
  $C_\mathrm{tot}=C_\mathrm{deg}+\mathrm{VOLL}(\mathrm{LPSP})\times
  E_\mathrm{unserved}$ across the sensitivity cases (rank~1 = lowest cost; the
  rightmost column is the mean rank). RB2(SoH) ranks first in every case.}
  \label{fig:ranking}
\end{figure}
